\subsection{复现步骤}

本仓库的设计原则是“唯一计算面”：一切数据解析、数值计算、优化求解、测试、绘图、\LaTeX{} 排版与打包
都在云端（Modal）执行，本地只负责编辑文本、提交任务、下载产物并核对 SHA-256。
\textbf{关于上表的时序说明。}该表由 \texttt{paper} 阶段在编译本 PDF 之前生成，因此它记录的是\emph{那一刻}的运行状态：
\texttt{paper} 一行显示为 \texttt{running}（其自身状态尚未写入），而 \texttt{qa}、\texttt{package}、\texttt{release}
三个阶段按依赖顺序排在 \texttt{paper} 之后，其行要么缺失、要么是前一次编译循环留下的状态与输出摘要。
这三行\emph{不应}被用作审计线索：交付物的权威记录是发布目录中的 \nolinkurl{release_manifest.json}
（逐文件 SHA-256）与支撑材料内的 \nolinkurl{REPRODUCE.md}、\nolinkurl{MANIFEST.json}，它们都在 \texttt{paper} 之后生成。
复现步骤如下：

\begin{enumerate}[leftmargin=1.8em,itemsep=0.15em]
  \item 安装 Modal 客户端并完成鉴权后，在仓库根目录执行 \texttt{python3 tools/cli.py provision}
        验证云端工具链（\TeX{} Live、中文字体、CP-SAT、HiGHS）。
  \item 数据接入与契约校验：\texttt{python3 tools/cli.py run ingest,validate,detect -{}-new-run -{}-size small}。
  \item 两个重阶段分离运行（各约 $40$ 分钟）：\texttt{run resolve -{}-size large -{}-spawn} 与
        \texttt{run resolve\_interval -{}-size large -{}-spawn}，随后 \texttt{wait <ID>} 或轮询 \texttt{status}。
  \item 其余科学阶段：\texttt{run pack,bench,sensitivity -{}-size large}，再 \texttt{run compare,results,figures,tables}。
  \item 门禁与论文：\texttt{run lint,test} 与 \texttt{run paper,qa,package,release}。
  \item 下载与核对：\texttt{download -{}-stage <name>} 会逐文件核对 SHA-256；
        \texttt{release -{}-version <tag>} 生成交付目录与发布清单。
  \item \textbf{不使用云端账号的复现路径}：支撑材料根目录的 \nolinkurl{tools/run_local.py} 调用同一批阶段函数，
        只把云端卷换成本地目录。安装 \texttt{pandas numpy openpyxl pyarrow networkx matplotlib ortools highspy} 后执行
        \texttt{python3 tools/run\_local.py -{}-quick}，数分钟内得到四个结果工作簿与全部图表；
        去掉 \texttt{-{}-quick} 则使用与本文相同量级的求解预算。排版阶段需另行安装 \TeX{} Live 与中文字体。
\end{enumerate}
各阶段的确切参数（求解器线程数、逐级时限）记录在仓库的运行手册中，也记录在每个阶段的 \nolinkurl{manifest.json} 内。

\subsection{可复现性的保证}

\begin{itemize}[leftmargin=1.6em,itemsep=0.15em]
  \item \textbf{固定输入}：题目 PDF 与附件在仓库中原样保存，\nolinkurl{ingest} 阶段记录其 SHA-256，任何代码不得改写。
  \item \textbf{固定环境}：容器镜像的依赖版本固定；每个阶段的 \nolinkurl{manifest.json} 记录 Python 版本、
        平台、CPU 数与全部包版本。
  \item \textbf{固定随机源}：所有随机性经统一的种子接口派生，求解器的随机种子固定。
  \item \textbf{固定热启动}：词典序求解的现任解提示保存在仓库的 \nolinkurl{configs/hints/} 中
        （附带来源 run、代码版本与目标向量）。由命题~\ref{prop:cut}，带提示重跑所得向量在词典序上不会劣于提示，
        因此本文结果可由仓库直接复现。提示只影响搜索起点，不改变可行域，也不改变任何一级的最优值。
  \item \textbf{数字的单一来源}：论文正文不允许手工键入结果数字；每个数值或来自阶段登记的宏，
        或来自 \nolinkurl{tables} 阶段生成的表格文件。若两个阶段登记了同名而不同值的数字，论文构建会报告冲突并由质检门禁拦截。
\end{itemize}

\subsection{已知的不确定性}

$\valQonePlans$ 个计划的完整模型采用多线程求解，单次求解耗时数十分钟，本文未做重复求解以验证位对位的确定性
（记为“\valQtwoDeterministic{}”）；微型实例上的确定性检验为 $\valBenchTinyDeterministic/\valBenchTinyInstances$ 通过。
需要强调：本文\emph{已证明最优}的结论（定理~\ref{thm:q2opt}、\ref{thm:q3}）由求解器的证明下界保证，
与是否复现出同一个最优解无关；\emph{未证明最优}的层级则本就以区间形式报告。
